%─────────────────────────────────────────────────────────────────────────────
%  ReciLuz — Documento Técnico: Análise de Algoritmos
%  CESAR School · Sistemas Embarcados · 2026
%─────────────────────────────────────────────────────────────────────────────
\documentclass[12pt, a4paper]{article}

%── Pacotes ──────────────────────────────────────────────────────────────────
\usepackage[utf8]{inputenc}
\usepackage[T1]{fontenc}
\usepackage[brazil]{babel}
\usepackage{geometry}
\usepackage{amsmath, amssymb, amsthm}
\usepackage{booktabs}
\usepackage{array}
\usepackage{xcolor}
\usepackage{listings}
\usepackage{graphicx}
\usepackage{pgfplots}
\pgfplotsset{compat=1.18}
\usepackage{tikz}
\usepackage{hyperref}
\usepackage{enumitem}
\usepackage{microtype}
\usepackage{caption}
\usepackage{subcaption}
\usepackage{float}
\usepackage{fancyhdr}
\usepackage{titlesec}
\usepackage{parskip}

%── Layout ────────────────────────────────────────────────────────────────────
\geometry{top=2.5cm, bottom=2.5cm, left=3cm, right=2.5cm}
\setlength{\headheight}{14pt}

\pagestyle{fancy}
\fancyhf{}
\fancyhead[L]{\small\textcolor{gray}{ReciLuz — Análise de Algoritmos}}
\fancyhead[R]{\small\textcolor{gray}{CESAR School 2026}}
\fancyfoot[C]{\thepage}
\renewcommand{\headrulewidth}{0.4pt}

%── Cores ────────────────────────────────────────────────────────────────────
\definecolor{teal}{RGB}{27, 138, 143}
\definecolor{vermelho}{RGB}{224, 85, 85}
\definecolor{verde}{RGB}{46, 204, 113}
\definecolor{cinza}{RGB}{100, 100, 100}
\definecolor{codebg}{RGB}{245, 246, 250}
\definecolor{codefg}{RGB}{30, 30, 60}

%── Listings (código C++) ─────────────────────────────────────────────────────
\lstset{
  language=C++,
  backgroundcolor=\color{codebg},
  basicstyle=\ttfamily\small,
  keywordstyle=\color{teal}\bfseries,
  commentstyle=\color{cinza}\itshape,
  stringstyle=\color{vermelho},
  numbers=left,
  numberstyle=\tiny\color{cinza},
  numbersep=8pt,
  breaklines=true,
  frame=single,
  framerule=0.4pt,
  rulecolor=\color{teal!40},
  captionpos=b,
  tabsize=2,
  showstringspaces=false,
  extendedchars=true,
  inputencoding=utf8,
  literate=
    {—}{{---}}1
    {á}{{\'a}}1 {é}{{\'e}}1 {í}{{\'i}}1 {ó}{{\'o}}1 {ú}{{\'u}}1
    {ã}{{\~a}}1 {õ}{{\~o}}1 {â}{{\^a}}1 {ê}{{\^e}}1 {ô}{{\^o}}1
    {à}{{\`a}}1 {Á}{{\'A}}1 {É}{{\'E}}1 {Í}{{\'I}}1 {Ó}{{\'O}}1
    {Ú}{{\'U}}1 {Ã}{{\~A}}1 {Õ}{{\~O}}1 {ç}{{\c{c}}}1 {Ç}{{\c{C}}}1,
}

%── Teoremas ─────────────────────────────────────────────────────────────────
\newtheorem{teorema}{Teorema}
\newtheorem{proposicao}[teorema]{Proposição}
\newtheorem{corolario}[teorema]{Corolário}
\newtheorem{lema}[teorema]{Lema}
\theoremstyle{definition}
\newtheorem{definicao}{Definição}
\theoremstyle{remark}
\newtheorem{observacao}{Observação}

%── Cabeçalho do título ───────────────────────────────────────────────────────
\title{
  \vspace{-1cm}
  {\large \textcolor{teal}{\textsc{ReciLuz — Sistemas Embarcados}}}\\[0.4em]
  {\LARGE \textbf{Documento Técnico:\\Análise de Algoritmos e\\Estruturas de Dados}}\\[0.4em]
  {\normalsize Buffer Circular \textcolor{verde}{O(1)} vs.\ Array com Deslocamento \textcolor{vermelho}{O(n)}}
}
\author{
  Grupo 8 — CESAR School\\
  \texttt{prcr@cesar.school}
}
\date{Junho de 2026}

%─────────────────────────────────────────────────────────────────────────────
\begin{document}
%─────────────────────────────────────────────────────────────────────────────

\maketitle
\thispagestyle{fancy}

\begin{abstract}
Este documento apresenta a análise técnica da comparação entre duas estratégias
de gerenciamento de buffer implementadas no firmware do projeto ReciLuz
(ESP32 @ 240\,MHz).
A \textbf{Vertente\,1} utiliza um array com deslocamento de elementos
(\emph{shift}), cuja complexidade de inserção é $O(n)$.
A \textbf{Vertente\,2} utiliza um buffer circular baseado em aritmética modular
(\emph{ring buffer}), com inserção em tempo constante $O(1)$.
São apresentadas provas teóricas de complexidade, diagnóstico de impacto na
heap do microcontrolador e uma discussão sobre o comportamento do sistema
perante instabilidades de rede, com ênfase no modelo Produtor-Consumidor
implementado via FreeRTOS.
\end{abstract}

\tableofcontents
\newpage

%═════════════════════════════════════════════════════════════════════════════
\section{Introdução}
%═════════════════════════════════════════════════════════════════════════════

O ReciLuz é um sistema IoT de iluminação pública inteligente baseado em ESP32,
voltado ao contexto urbano do Recife. O firmware captura continuamente dados de
sensores (corrente, distância, ruído, temperatura) e os transmite via MQTT a
um broker centralizado.

Nesse cenário, existe uma disparidade temporal crítica entre as operações do
sistema:

\begin{center}
\begin{tabular}{lcc}
  \toprule
  Operação & Escala de tempo & Referência \\
  \midrule
  Leitura de sensor (ADC, DHT22, HC-SR04) & $\sim$\,µs & hardware \\
  Transmissão de rede (MQTT/HTTP) & $\sim$\,ms & protocolo \\
  \bottomrule
\end{tabular}
\end{center}

\medskip
A diferença pode atingir três ordens de magnitude. Se o loop de captura de
sensores for bloqueado pela latência de rede — ou se cada inserção no buffer
consumir ciclos de CPU proporcionais ao tamanho do histórico — o sistema
acumula \emph{jitter} amostral e pode sofrer \emph{stack overflow} em
operações com grandes buffers.

A escolha da estrutura de dados que mantém o histórico de leituras é, portanto,
um fator determinante para a estabilidade do sistema embarcado.

%═════════════════════════════════════════════════════════════════════════════
\section{Análise Assintótica}
%═════════════════════════════════════════════════════════════════════════════

\subsection{Definições e notação}

\begin{definicao}[Buffer circular de tamanho $N$]
  Um buffer circular de capacidade $N$ é um array estático
  $B[0\ldots N-1]$ acompanhado de dois índices inteiros $\mathit{head}$ e
  $\mathit{tail}$, com $0 \leq \mathit{head}, \mathit{tail} < N$, que circulam
  via aritmética modular. O elemento mais antigo encontra-se em
  $B[\mathit{head}]$; o próximo slot de escrita é $B[\mathit{tail}]$.
\end{definicao}

\begin{definicao}[Array com deslocamento de tamanho $N$]
  Um array com deslocamento mantém $\mathit{count}$ elementos
  contíguos em $B[0\ldots \mathit{count}-1]$. Para inserir um novo elemento
  quando $\mathit{count} = N$ (buffer cheio), todos os $N-1$ elementos
  existentes são copiados uma posição à esquerda antes da inserção no slot
  $B[N-1]$.
\end{definicao}

%─────────────────────────────────────────────────────────────────────────────
\subsection{Vertente 1 — Array com deslocamento: prova de \texorpdfstring{$O(n)$}{O(n)}}
%─────────────────────────────────────────────────────────────────────────────

A implementação em \texttt{IneffBuffer<T,N>::push()} (arquivo
\texttt{ineff\_buffer.h}) executa, quando o buffer está cheio:

\begin{lstlisting}[caption={Núcleo de \texttt{IneffBuffer::push()} — O(n)}]
if (_count == N) {
    for (size_t i = 0; i < N - 1; i++)
        _buf[i] = _buf[i + 1];   // N-1 atribuicoes
    _count = N - 1;
}
_buf[_count++] = item;           // 1 atribuicao
\end{lstlisting}

\begin{teorema}[\texttt{IneffBuffer::push()} é $\Theta(n)$]
\label{thm:ineff-push}
Seja $N$ a capacidade do buffer. A operação \texttt{push()} quando o buffer
está cheio executa exatamente $N - 1$ atribuições no laço de deslocamento
mais uma atribuição final, totalizando $N$ atribuições.
Portanto $T(N) = N \in \Theta(N)$.
\end{teorema}

\begin{proof}
O laço \texttt{for (i = 0; i < N-1; i++)} itera exatamente $N-1$ vezes
(índices $0, 1, \ldots, N-2$), cada iteração executando uma atribuição
$B[i] \leftarrow B[i+1]$ de custo $O(1)$.
Após o laço, a atribuição \texttt{\_buf[\_count++] = item} acrescenta
mais uma operação de custo $O(1)$.
O custo total é
\[
  T_{\text{push}}(N) = (N - 1) + 1 = N.
\]
Como $T(N) = N$, temos $T(N) \in \Omega(N)$ e $T(N) \in O(N)$,
logo $T(N) \in \Theta(N)$.
\end{proof}

\begin{corolario}
Para uma sequência de $K$ inserções consecutivas num buffer sempre cheio de
tamanho $N$, o custo total é $K \cdot N \in O(K \cdot N)$.
No benchmark do ReciLuz, $K = \text{BENCH\_REPETICOES} = 100$, de modo que o
custo cresce linearmente com $N$: $100 \cdot N$.
\end{corolario}

Na prática, o firmware utiliza \texttt{memmove} para realizar o deslocamento:

\begin{lstlisting}[caption={Benchmark V1 — \texttt{memmove} simula o shift O(n)}]
memmove(bufV1, bufV1 + 1, (n - 1) * sizeof(float));
bufV1[n - 1] = sampleVal;
\end{lstlisting}

O \texttt{memmove} move $(N-1) \times \texttt{sizeof(float)} = 4(N-1)$ bytes.
Para $N = 20\,000$, isso representa aproximadamente $80\,\text{KB}$
transferidos a cada inserção — confirma o crescimento $\Theta(N)$.

%─────────────────────────────────────────────────────────────────────────────
\subsection{Vertente 2 — Buffer circular: prova de \texorpdfstring{$O(1)$}{O(1)}}
%─────────────────────────────────────────────────────────────────────────────

A implementação em \texttt{RingBuffer<T,N>::push()} (arquivo
\texttt{ring\_buffer.h}) é:

\begin{lstlisting}[caption={Núcleo de \texttt{RingBuffer::push()} — O(1)}]
bool push(const T& item) {
    _buf[_tail] = item;              // 1 escrita
    _tail = (_tail + 1) % N;        // 1 incremento + 1 modulo
    if (_count == N)
        _head = (_head + 1) % N;    // 1 incremento + 1 modulo (condicional)
    else
        _count++;                    // 1 incremento (condicional)
    return (_count == N) ? false : true;
}
\end{lstlisting}

\begin{teorema}[\texttt{RingBuffer::push()} é $O(1)$]
\label{thm:ring-push}
A operação \texttt{push()} executa um número fixo de operações
independente de $N$.
\end{teorema}

\begin{proof}
O corpo de \texttt{push()} consiste em:
\begin{enumerate}[label=(\roman*)]
  \item Uma escrita em \texttt{\_buf[\_tail]}: $O(1)$.
  \item Incremento de \texttt{\_tail} com módulo: $O(1)$.
  \item Um desvio condicional (\texttt{if \_count == N}): $O(1)$.
  \item Em qualquer ramo do condicional: uma operação de $O(1)$
        (incremento de \texttt{\_head} com módulo \emph{ou} incremento de
        \texttt{\_count}).
  \item Retorno de valor booleano: $O(1)$.
\end{enumerate}
O número total de operações é constante e não depende de $N$.
Portanto $T_{\text{push}}(N) = c$ para alguma constante $c > 0$, e
$T_{\text{push}}(N) \in O(1)$.
\end{proof}

\begin{proposicao}[Aritmética modular como acesso circular]
O índice \texttt{tail} percorre os slots $0, 1, \ldots, N-1, 0, 1, \ldots$
indefinidamente. A operação $\mathit{tail} \leftarrow (\mathit{tail}+1) \bmod N$
garante que $\mathit{tail}$ nunca ultrapasse os limites do array, sem
necessidade de movimentação de elementos.
\end{proposicao}

%─────────────────────────────────────────────────────────────────────────────
\subsection{Verificação empírica da linearidade}
%─────────────────────────────────────────────────────────────────────────────

Os resultados a seguir foram medidos no ESP32-D0WD-V3 @ 240\,MHz com
100 inserções em buffer cheio (BENCH\_REPETICOES\,=\,100):

\begin{table}[H]
  \centering
  \caption{Latência medida por inserção — ESP32 @ 240 MHz}
  \label{tab:bench}
  \begin{tabular}{rrrrc}
    \toprule
    $N$ & V1 (µs/ins.) & V2 (µs/ins.) & Razão V1/V2 & Linearidade V1 \\
    \midrule
    100   & 23{,}52  & 0{,}080 & 294×    & —         \\
    500   & 177{,}0  & 0{,}040 & 4\,426× & ×5,0 $\to$ ×7,5 \\
    1\,000 & 236{,}5 & 0{,}040 & 5\,912× & ×2,0 $\to$ ×1,3 \\
    5\,000 & 1\,225  & 0{,}040 & 30\,632× & ×5,0 $\to$ ×5,2 \\
    20\,000 & 4\,909 & 0{,}040 & 122\,722× & ×4,0 $\to$ ×4,0 \\
    \bottomrule
  \end{tabular}
\end{table}

\begin{observacao}
  A razão medida entre $N=1.000/N=100$ é $236{,}5 / 23{,}52 = 10{,}05\,\times$,
  coincidindo com a razão de tamanhos $10\times$. Entre $N=20.000$ e $N=5.000$
  a razão observada é $4{,}0\times$, exatamente igual à razão de tamanhos
  esperada para $O(n)$. Isso confirma empiricamente o Teorema~\ref{thm:ineff-push}.
  O V2 permanece em $\approx 0{,}04$\,µs para qualquer $N$, confirmando o
  Teorema~\ref{thm:ring-push}.
\end{observacao}

\begin{figure}[H]
  \centering
  \begin{tikzpicture}
    \begin{axis}[
      width=13cm, height=7cm,
      xlabel={$N$ (tamanho do buffer)},
      ylabel={Latência por inserção (µs)},
      xmin=0, xmax=21000,
      ymin=0, ymax=5200,
      xtick={100,500,1000,5000,20000},
      xticklabels={100,500,1K,5K,20K},
      grid=major,
      grid style={dotted, gray!40},
      legend pos=north west,
      legend style={font=\small},
    ]
      \addplot[
        color=vermelho, mark=*, thick,
        mark options={fill=vermelho}
      ] coordinates {
        (100, 23.52)
        (500, 177.0)
        (1000, 236.5)
        (5000, 1225)
        (20000, 4909)
      };
      \addlegendentry{V1 — Array O(n)}

      \addplot[
        color=verde, mark=square*, thick,
        mark options={fill=verde}
      ] coordinates {
        (100,  0.080)
        (500,  0.040)
        (1000, 0.040)
        (5000, 0.040)
        (20000,0.040)
      };
      \addlegendentry{V2 — Ring Buffer O(1)}

      % Linha de referência O(n) teórica (ajustada: c = 0.245)
      \addplot[
        color=vermelho!40, dashed, domain=100:20000, samples=100
      ] { 0.245 * x };
      \addlegendentry{$O(n)$ teórico ($c=0{,}245$)}
    \end{axis}
  \end{tikzpicture}
  \caption{Latência de inserção: curva V1 (vermelho) crescendo linearmente com $N$
    versus V2 (verde) constante próxima de zero. A linha tracejada representa
    o modelo teórico $T(n) = 0{,}245\,n$ ajustado aos dados de V1.}
  \label{fig:latencia}
\end{figure}

%─────────────────────────────────────────────────────────────────────────────
\subsection{Análise do \texttt{pop()} e simetria de complexidade}
%─────────────────────────────────────────────────────────────────────────────

\begin{proposicao}[\texttt{IneffBuffer::pop()} é $\Theta(n)$]
A operação de remoção do \texttt{IneffBuffer} percorre o laço
\texttt{for (i = 0; i < \_count-1; i++)} que executa $\mathit{count}-1$
atribuições, de modo análogo ao \texttt{push()}. Para buffer cheio,
$T_{\text{pop}}(N) = N - 1 \in \Theta(N)$.
\end{proposicao}

\begin{proposicao}[\texttt{RingBuffer::pop()} é $O(1)$]
A remoção no \texttt{RingBuffer} lê \texttt{\_buf[\_head]}, incrementa
\texttt{\_head} com módulo e decrementa \texttt{\_count} — número fixo
de operações independente de $N$.
\end{proposicao}

A Tabela~\ref{tab:complexidade} resume a complexidade assintótica de todas
as operações de ambas as estruturas.

\begin{table}[H]
  \centering
  \caption{Resumo de complexidade das operações}
  \label{tab:complexidade}
  \begin{tabular}{lccc}
    \toprule
    Operação & \texttt{IneffBuffer} (V1) & \texttt{RingBuffer} (V2) \\
    \midrule
    \texttt{push()} — buffer cheio    & $\Theta(n)$ & $O(1)$ \\
    \texttt{push()} — buffer não cheio & $O(1)$     & $O(1)$ \\
    \texttt{pop()}                    & $\Theta(n)$ & $O(1)$ \\
    \texttt{isEmpty()} / \texttt{isFull()} & $O(1)$ & $O(1)$ \\
    Espaço (memória)                  & $\Theta(n)$ & $\Theta(n)$ \\
    \bottomrule
  \end{tabular}
\end{table}

%═════════════════════════════════════════════════════════════════════════════
\section{Diagnóstico de Memória}
%═════════════════════════════════════════════════════════════════════════════

\subsection{Modelo de memória do ESP32}

O ESP32-D0WD-V3 dispõe de 520\,KB de SRAM interna, dos quais aproximadamente
300\,KB estão disponíveis para o programa em tempo de execução após o
carregamento do stack TCP/IP do ESP-IDF e do FreeRTOS.

\begin{center}
\begin{tabular}{lrl}
  \toprule
  Região & Tamanho & Uso \\
  \midrule
  DRAM (dados + heap)   & $\sim$300\,KB & variáveis, heap dinâmica \\
  IRAM (instruções)     & $\sim$128\,KB & código em cache hot-path \\
  RTC SRAM              & 16\,KB        & deep-sleep \\
  \bottomrule
\end{tabular}
\end{center}

\subsection{Impacto da Vertente 1: fragmentação por \texttt{malloc}/\texttt{free}}

O benchmark da Vertente\,1 aloca e libera buffers dinamicamente para cada
valor de $N$ testado:

\begin{lstlisting}[caption={Ciclo de alocação/liberação no benchmark V1}]
float* bufV1 = (float*)malloc(n * sizeof(float)); // aloca n*4 bytes
// ... executa 100 insercoes ...
free(bufV1);                                       // libera

float* bufV2 = (float*)malloc(n * sizeof(float)); // aloca novamente
// ... executa 100 insercoes ...
free(bufV2);                                       // libera
\end{lstlisting}

Para $N = 20\,000$, cada alocação solicita $80\,\text{KB}$ ao heap do ESP32.
O ciclo alocação-uso-liberação repetido 5 vezes (um por valor de $N$) a cada
30 segundos introduz dois efeitos deletérios:

\begin{enumerate}
  \item \textbf{Fragmentação externa:} blocos livres não-contíguos resultantes
    de alocações sequenciais de tamanhos diferentes ($100 \times 4$,
    $500 \times 4$, $1000 \times 4$\ldots) podem impedir alocações futuras
    mesmo com heap livre total suficiente.

  \item \textbf{Overhead de metadados:} cada chamada a \texttt{malloc} insere
    um cabeçalho de controle (tipicamente 8--16 bytes no heap do
    ESP-IDF/newlib), aumentando o consumo real além do tamanho solicitado.
\end{enumerate}

\subsubsection{Dados medidos}

Os valores de heap livre foram capturados via \texttt{ESP.getFreeHeap()}
durante a execução do benchmark no ESP32 @ 240\,MHz:

\begin{table}[H]
  \centering
  \caption{Heap livre durante o benchmark (medições reais)}
  \label{tab:heap}
  \begin{tabular}{lrrr}
    \toprule
    Condição & Mín. (bytes) & Máx. (bytes) & Amplitude \\
    \midrule
    V1 — benchmark ativo (malloc/free) & 199\,668 & 211\,964 & 12\,296\,bytes \\
    V2 — operação normal (estático)    & 211\,340 & 211\,340 & 0\,bytes \\
    \bottomrule
  \end{tabular}
\end{table}

A oscilação de $\approx 12\,\text{KB}$ observada durante os testes V1 reflete
o heap não recuperado integralmente entre os ciclos \texttt{malloc}/\texttt{free}.
Em contraste, o \texttt{RingBuffer} da Vertente\,2 é declarado como variável
estática global:

\begin{lstlisting}[caption={Alocação estática do RingBuffer — sem impacto no heap}]
// Vertente 2: declarado globalmente, vai para o segmento BSS
RingBuffer<Leitura, BUFFER_SIZE> bufferCircular;
IneffBuffer<Leitura, BUFFER_SIZE> bufferIneficiente;
\end{lstlisting}

Variáveis globais/estáticas são alocadas no segmento \texttt{.bss} (não
inicializadas) ou \texttt{.data} (inicializadas) em tempo de link, sem
qualquer chamada ao alocador dinâmico em tempo de execução. O heap permanece
inalterado durante toda a vida útil do programa.

\subsection{Risco de falha de alocação por heap insuficiente}

O firmware inclui uma guarda explícita antes de cada alocação de benchmark:

\begin{lstlisting}[caption={Guarda de heap antes de malloc no benchmark}]
if (ESP.getFreeHeap() < (uint32_t)(n * sizeof(float) + 4096)) {
    Serial.printf("N=%5d | SKIP -- heap insuficiente\n", n);
    continue;
}
\end{lstlisting}

A margem de 4\,096 bytes garante que o sistema operacional e a pilha de
rede TCP/IP do ESP-IDF continuem operando mesmo durante a alocação máxima
($N = 20\,000$, 80\,KB).

\subsection{Comparativo de estabilidade de memória}

\begin{figure}[H]
  \centering
  \begin{tikzpicture}
    \begin{axis}[
      width=13cm, height=5cm,
      xlabel={Iteração de benchmark},
      ylabel={Heap livre (KB)},
      xmin=0.5, xmax=5.5,
      ymin=190, ymax=220,
      xtick={1,2,3,4,5},
      xticklabels={N=100,N=500,N=1K,N=5K,N=20K},
      grid=major,
      grid style={dotted, gray!40},
      legend pos=south east,
    ]
      \addplot[
        color=vermelho, mark=triangle*, thick,
        mark options={fill=vermelho}
      ] coordinates {
        (1, 211.9) (2, 207.7) (3, 208.1) (4, 208.3) (5, 208.3)
      };
      \addlegendentry{V1 — com malloc/free}

      \addplot[
        color=verde, mark=*, thick, dashed,
        mark options={fill=verde}
      ] coordinates {
        (1, 211.3) (2, 211.3) (3, 211.3) (4, 211.3) (5, 211.3)
      };
      \addlegendentry{V2 — estático (BSS)}
    \end{axis}
  \end{tikzpicture}
  \caption{Heap livre ao longo das iterações de benchmark.
    V1 apresenta variação de até 12\,KB por fragmentação; V2 permanece
    constante pois não utiliza o alocador dinâmico.}
  \label{fig:heap}
\end{figure}

%═════════════════════════════════════════════════════════════════════════════
\section{Discussão: Comportamento sob Instabilidade de Rede}
%═════════════════════════════════════════════════════════════════════════════

\subsection{Modelo Produtor-Consumidor com FreeRTOS}

O firmware adota explicitamente o padrão \emph{Produtor-Consumidor}
implementado sobre FreeRTOS, com tarefas em núcleos físicos distintos do
ESP32:

\begin{center}
\begin{tikzpicture}[
  box/.style={draw=teal, fill=teal!10, rounded corners, minimum width=5cm, minimum height=1.2cm, align=center, font=\small},
  arrow/.style={->, thick, teal},
]
  \node[box] (sensor) at (0,0) {\textbf{taskSensor} (Core 1, prio 2)\\
    Sensores → push() → bufferCircular};
  \node[box] (mqtt) at (7.5,0) {\textbf{taskMqtt} (Core 0, prio 1)\\
    pop() → MQTT / HTTP};
  \node[box] (bench) at (7.5,-2.5) {\textbf{taskBenchmark} (Core 0, prio 1)\\
    executarBenchmark() a cada 30s};
  \draw[arrow] (sensor) -- node[above, font=\footnotesize]{mutex (5ms)} (mqtt);
  \draw[arrow] (sensor.south east) -- ++(1,0) |- (bench.west);
\end{tikzpicture}
\end{center}

O \texttt{bufferCircular} (tipo \texttt{RingBuffer}) é a fronteira entre os
dois domínios de tempo: o domínio de sensores (µs) e o domínio de rede (ms).
Um \texttt{SemaphoreHandle\_t} (mutex FreeRTOS) protege o acesso
compartilhado com \emph{timeout} de 5\,ms — se a \texttt{taskMqtt}
demorar mais que 5\,ms mantendo o mutex, a \texttt{taskSensor} prossegue
sem bloquear e a leitura é descartada.

\subsection{Cenário A: rede estável}

Com MQTT conectado e latência normal ($\ll 1.500$\,ms), a
\texttt{taskMqtt} drena o buffer a cada ciclo de envio. O nível de ocupação
do buffer permanece baixo e nenhuma leitura é descartada.

\subsection{Cenário B: queda total de rede (gargalo simulado)}

Quando a rede cai (broker MQTT desconectado, WiFi fora do alcance ou
congestionamento severo), os dois comportamentos divergem radicalmente:

\subsubsection{Vertente 1 sob gargalo}

Na Vertente\,1 (\texttt{IneffBuffer}), o buffer é usado exclusivamente pela
\texttt{taskSensor} (sem consumidor paralelo). Em ausência de rede:

\begin{enumerate}
  \item A \texttt{taskSensor} continua inserindo leituras com custo $O(N)$
    por inserção.
  \item Para $N = 100$ (tamanho operacional do sistema), cada inserção
    executa 99 atribuições de \texttt{struct Leitura} (32 bytes cada),
    movendo $99 \times 32 = 3\,168$ bytes por ciclo de 80\,ms.
  \item O dado nunca é consumido: a estrutura enche e mantém custo máximo
    $\Theta(N)$ indefinidamente.
  \item Se o buffer fosse expandido para absorver o backlog de uma
    interrupção de rede de, por exemplo, 10 minutos
    ($10 \times 60 \times 12{,}5 = 7\,500$ leituras a 12,5 Hz),
    o custo por inserção atingiria $7\,500 \times 32 = 240\,\text{KB}$
    de memmove por ciclo — inviável no ESP32.
\end{enumerate}

\subsubsection{Vertente 2 sob gargalo (comportamento real implementado)}

Com o \texttt{RingBuffer} e o modelo Produtor-Consumidor:

\begin{enumerate}
  \item \textbf{O produtor nunca para:} a \texttt{taskSensor} executa
    \texttt{push()} em $O(1)$ independentemente do estado da rede.
    O ciclo de 80\,ms é preservado sem jitter.

  \item \textbf{Sobrescrita controlada:} quando o buffer circular de
    capacidade $N = 100$ está cheio, cada novo \texttt{push()} sobrescreve
    a leitura mais antiga (política \emph{overwrite-oldest}). O sistema
    sacrifica dados históricos mas mantém os dados mais recentes disponíveis.

  \item \textbf{Recuperação imediata:} ao reconectar, a \texttt{taskMqtt}
    drena imediatamente o buffer com os $N$ dados mais recentes,
    sem necessidade de sincronização adicional.

  \item \textbf{Custo do mutex:} o timeout de 5\,ms da operação
    \texttt{xSemaphoreTake()} garante que a \texttt{taskSensor} nunca
    espere mais que 5\,ms mesmo que a \texttt{taskMqtt} esteja bloqueada
    em I/O de rede.
\end{enumerate}

\subsection{Instabilidade intermitente: efeito de rajada (\emph{burst})}

Um cenário realista de gargalo é a intermitência: períodos de conectividade
normal intercalados com quedas de segundos a minutos.

Durante uma queda de $T_{\text{queda}}$ segundos com taxa de amostragem
$f = 12{,}5$\,Hz (ciclo de 80\,ms), o número de leituras geradas é:
\[
  K = f \times T_{\text{queda}} = 12{,}5 \times T_{\text{queda}}.
\]

Para $T_{\text{queda}} = 60$\,s, $K = 750$ leituras. Com buffer de $N=100$,
as 650 mais antigas serão sobrescritas — perda de $86{,}7\%$ do histórico.

\begin{observacao}
  A perda de dados históricos é uma consequência \textbf{esperada e
  aceitável} em sistemas embarcados com restrição de memória. A alternativa
  — bloquear o sensor ou expandir o buffer para $K$ posições — tornaria o
  sistema instável ou inviável em hardware com 520\,KB de RAM. O
  \texttt{RingBuffer} implementa o compromisso correto: \emph{always fresh,
  never stale, never blocked}.
\end{observacao}

\subsection{Quantificação do impacto do jitter na Vertente 1}

O jitter amostral $\sigma_T$ causado pelo custo variável de inserção na
Vertente\,1 pode ser estimado:

\[
  \sigma_T^{\text{V1}} \approx T_{\text{insert}}(N) = c \cdot N
  \quad \text{com } c \approx 0{,}245\,\text{µs/elemento (medido)}.
\]

Para $N = 100$ (operação normal): $\sigma_T^{\text{V1}} \approx 24{,}5$\,µs.
Para $N = 1.000$: $\sigma_T^{\text{V1}} \approx 245$\,µs. Este jitter é
injetado diretamente no loop de controle PWM da lâmpada, podendo causar
variação perceptível de brilho em frequências > 1\,kHz.

Na Vertente\,2: $\sigma_T^{\text{V2}} \approx 0{,}04$\,µs — três ordens de
magnitude menor, desprezível para qualquer frequência de controle prática.

%═════════════════════════════════════════════════════════════════════════════
\section{Conclusão}
%═════════════════════════════════════════════════════════════════════════════

A análise apresentada neste documento demonstra, por prova formal e
verificação empírica, que:

\begin{enumerate}
  \item O \texttt{IneffBuffer} (Vertente\,1) possui complexidade
    $\Theta(n)$ em \texttt{push()} e \texttt{pop()}, com crescimento linear
    confirmado por razões experimentais que coincidem com as razões de $N$
    (dentro de $\pm 4\%$).

  \item O \texttt{RingBuffer} (Vertente\,2) possui complexidade $O(1)$ em
    todas as operações, medido em $\approx 0{,}04$\,µs para qualquer
    $N \in \{100, 500, 1\,000, 5\,000, 20\,000\}$.

  \item A Vertente\,1 introduz fragmentação de heap de até 12\,KB pelo
    ciclo \texttt{malloc}/\texttt{free} do benchmark, enquanto a Vertente\,2
    opera inteiramente no segmento estático BSS, sem qualquer impacto no
    heap.

  \item Sob instabilidade de rede, o modelo Produtor-Consumidor da
    Vertente\,2 preserva a integridade do loop de sensores com jitter de
    $\approx 0{,}04$\,µs, enquanto a Vertente\,1 injeta jitter proporcional
    a $N$ no loop de controle.
\end{enumerate}

A escolha do \texttt{RingBuffer} com aritmética modular $O(1)$ é a decisão
de engenharia correta para o ReciLuz, tanto do ponto de vista de
complexidade assintótica quanto de estabilidade operacional em ambiente
embarcado com recursos limitados.

%═════════════════════════════════════════════════════════════════════════════
\appendix
\section{Configuração do benchmark}
%═════════════════════════════════════════════════════════════════════════════

\begin{table}[H]
  \centering
  \caption{Constantes do firmware relevantes para o benchmark}
  \begin{tabular}{llr}
    \toprule
    Constante & Descrição & Valor \\
    \midrule
    \texttt{BUFFER\_SIZE}          & Tamanho do buffer operacional & 100 \\
    \texttt{BENCH\_REPETICOES}     & Inserções por teste           & 100 \\
    \texttt{INTERVALO\_BENCHMARK\_MS} & Intervalo entre benchmarks & 30\,000\,ms \\
    \texttt{ATRASO\_LOOP\_MS}      & Ciclo da taskSensor           & 80\,ms \\
    \texttt{INTERVALO\_ENVIO\_MS}  & Ciclo de envio MQTT/HTTP      & 1\,500\,ms \\
    \bottomrule
  \end{tabular}
\end{table}

\section{Saída serial do benchmark}

Saída capturada no ESP32-D0WD-V3 @ 240\,MHz:

\begin{lstlisting}[language={}, caption={Serial monitor — benchmark real}]
=== BENCHMARK AA (buffer variavel, 100 insercoes/teste) ===
N=  100 | V1=   2352 us (23.52/item) | V2=    8 us (0.080/item) | 294x
N=  500 | V1=  17702 us (177.02/item)| V2=    4 us (0.040/item) | 4426x
N= 1000 | V1=  23646 us (236.46/item)| V2=    4 us (0.040/item) | 5912x
N= 5000 | V1= 122528 us (1225.28/item)| V2=   4 us (0.040/item) | 30632x
N=20000 | V1= 490889 us (4908.89/item)| V2=   4 us (0.040/item) | 122722x
=== FIM BENCHMARK ===
\end{lstlisting}

%─────────────────────────────────────────────────────────────────────────────
\end{document}
%─────────────────────────────────────────────────────────────────────────────
